\documentclass[11pt,a4paper]{article}
\usepackage[utf8]{inputenc}
\usepackage[T1]{fontenc}
\usepackage{lmodern}
\usepackage[french]{babel}
\usepackage{amsmath,amssymb,mathtools}
\usepackage{icomma}
\usepackage{xcolor}
\usepackage[most]{tcolorbox}
\usepackage{enumitem}
\usepackage[a4paper,margin=2.2cm,headheight=26pt,headsep=10pt]{geometry}
\usepackage{fancyhdr}
\usepackage[hidelinks]{hyperref}
\definecolor{primary}{RGB}{222,49,99}
\definecolor{primarydark}{RGB}{150,22,55}
\definecolor{excol}{RGB}{0,135,90}
\tcbset{boxbase/.style={enhanced,breakable,boxrule=0.9pt,arc=2.5pt,
  left=3mm,right=3mm,top=2.3mm,bottom=2.3mm,fonttitle=\bfseries,coltitle=white}}
\newtcolorbox[auto counter]{corrige}[1][]{boxbase,colback=excol!4,colframe=excol,
  title={\textsc{Corrigé}~\thetcbcounter\ifx\relax#1\relax\relax\else\textmd{ --- \itshape #1}\fi}}
\newcommand{\R}{\mathbb{R}}
\newcommand{\ds}{\displaystyle}
\pagestyle{fancy}\fancyhf{}
\fancyhead[L]{\small\textcolor{primarydark}{\textbf{Thème 3} — Factorisation et résolution}}
\fancyhead[R]{\small\textcolor{primarydark}{\textsc{Corrigé — Moyen}}}
\fancyfoot[C]{\small\thepage}
\renewcommand{\headrulewidth}{0.8pt}
\begin{document}

\begin{center}
{\LARGE\bfseries\color{primarydark} Niveau Moyen — Corrigé}\\[2pt]
{\color{primary}\rule{0.45\textwidth}{1.2pt}}
\end{center}
\vspace{1mm}

\begin{corrige}[mise en évidence]
\begin{align*}
\text{a) }&5\cdot 2^{2x}-3a\cdot 2^{2x}\cdot 2=2^{2x}(5-6a); &
\text{b) }&e^{x}(e^{x}+1);\\
\text{c) }&(x-1)\bigl[(2x+3)-(x-1)\bigr]=(x-1)(x+4); &
\text{d) }&2x^{2}(3x^{2}-2x+1).
\end{align*}
\end{corrige}

\begin{corrige}[produits remarquables composés]
\begin{align*}
\text{a) }&(e^{x}+1)^{2};\\
\text{b) }&e^{4x}-1=(e^{2x}-1)(e^{2x}+1)=(e^{x}-1)(e^{x}+1)(e^{2x}+1);\\
\text{c) }&5x^{2}-9=(\sqrt5\,x)^{2}-3^{2}=(\sqrt5\,x-3)(\sqrt5\,x+3);\\
\text{d) }&(\ln x^{2})^{2}-3=(\ln x^{2}-\sqrt3)(\ln x^{2}+\sqrt3)=(2\ln x-\sqrt3)(2\ln x+\sqrt3).
\end{align*}
\end{corrige}

\begin{corrige}[trinôme puis résolution]
\begin{itemize}[leftmargin=1.4em,topsep=1pt,itemsep=1pt]
  \item a) $\Delta=9-8=1$ ; $x=\tfrac{3\pm1}{4}$ donc $x=1$ ou $x=\tfrac12$. Factorisation $(x-1)(2x-1)$, $S=\{\tfrac12,1\}$.
  \item b) $\Delta=12-12=0$ ; racine double $x=\tfrac{2\sqrt3}{6}=\tfrac{\sqrt3}{3}$. Ici $3x^{2}-2\sqrt3x+1=(\sqrt3x-1)^{2}$, $S=\{\tfrac{\sqrt3}{3}\}$.
  \item c) $\Delta=1+24=25$ ; $x=\tfrac{1\pm5}{2}$ donc $x=3$ ou $x=-2$, $S=\{-2,3\}$.
\end{itemize}
\end{corrige}

\begin{corrige}[Horner]
\textbf{a)} $P(1)=1-2-1+2=0$. Horner sur $(1,-2,-1,2)$ donne $Q(x)=x^{2}-x-2=(x-2)(x+1)$, donc
\[ P(x)=(x-1)(x-2)(x+1). \]
\textbf{b)} $P(1)=1-1+2-2=0$. Horner sur $(1,-1,2,-2)$ donne $Q(x)=x^{2}+2$ (sans racine réelle), donc
\[ P(x)=(x-1)(x^{2}+2). \]
\end{corrige}

\begin{corrige}[fraction rationnelle]
\textbf{a)} $\dfrac{x^{2}-1}{x^{2}-3x+2}=\dfrac{(x-1)(x+1)}{(x-1)(x-2)}=\dfrac{x+1}{x-2}$, avec C.E. $x\neq1$ et $x\neq2$.\\
\textbf{b)} $\dfrac{x^{2}-4}{x+2}=\dfrac{(x-2)(x+2)}{x+2}=x-2$, avec C.E. $x\neq-2$.
\end{corrige}

\end{document}
